\chapter{Microdermabrasion}
\index{Microdermabrasion}

\section*{Definition and Overview}
Microdermabrasion is a minimally invasive, non-surgical dermatological procedure designed to exfoliate and remove the outermost layer of dead skin cells, specifically the stratum corneum. The procedure is performed using a specialised mechanical handpiece that either projects fine abrasive crystals (typically aluminium oxide or sodium bicarbonate) onto the skin and vacuums them away, or utilises a solid diamond-tipped wand combined with adjustable vacuum suction. By removing the superficial epidermal layer, microdermabrasion stimulates cellular turnover, enhances local microcirculation, and promotes collagen synthesis in the upper dermis. It is widely performed as a low-risk, clinical intervention to improve superficial skin texture, reduce fine rhytids, lighten mild hyperpigmentation, and treat superficial acne scarring.

\section*{Epidemiology}
Microdermabrasion is consistently ranked among the top five most frequently performed non-surgical cosmetic procedures globally. It is popular across a wide age range, with the peak demographic being adults aged 25 to 55 years seeking facial rejuvenation or correction of early photoageing. The treatment is equally popular among women and men. Unlike deeper skin resurfacing modalities (such as deep chemical peels or ablative lasers), microdermabrasion is safe and effective for all skin phenotypes on the Fitzpatrick scale (Types I through VI). This is because the mechanical exfoliation is superficial and does not generate thermal damage, resulting in an extremely low incidence of post-inflammatory hyperpigmentation in darker skin tones.

\section*{Aetiology and Pathophysiology}
The cosmetic and physiological benefits of microdermabrasion are achieved through mechanical disruption of the epidermal barrier and cellular stimulation.
\begin{itemize}
    \item \textbf{Mechanical exfoliation:} The abrasive action of the crystals or diamond tip physically detaches and removes desquamating corneocytes from the stratum corneum, instantly smoothing the skin's surface.
    \item \textbf{Vacuum-induced microcirculation:} The negative pressure generated by the vacuum suction increases local blood flow (causing transient hyperaemia) and enhances lymphatic drainage. This increased perfusion delivers oxygen and essential nutrients to the basal layers of the epidermis.
    \item \textbf{Epidermal regeneration:} Disruption of the stratum corneum stimulates the mitotic activity of basal keratinocytes, accelerating the epidermal transit time and leading to a fresher, more compact epidermal layer.
    \item \textbf{Collagen remodeling:} While the abrasion is superficial, the mechanical stress and increased blood flow stimulate dermal fibroblasts, leading to a mild upregulation of Type I collagen and elastin fibres in the papillary dermis over a series of treatments.
    \item \textbf{Enhanced transdermal permeability:} Removal of the stratum corneum barrier temporarily increases skin permeability, significantly enhancing the absorption and efficacy of topically applied dermatological agents (such as hyaluronic acid, vitamin C, or light chemical peels) administered immediately post-procedure.
\end{itemize}

\section*{Clinical Features}
Microdermabrasion is indicated for several superficial aesthetic concerns, and presents with predictable post-treatment skin changes:
\begin{itemize}
    \item \textbf{Indications for treatment:}
        \begin{itemize}
            \item Mild facial rhytids (fine lines) and early signs of sun damage.
            \item Superficial dyschromia, including solar lentigines (age spots), dull skin tone, and mild melasma.
            \item \textbf{Acne-related concerns:} Comedonal acne, enlarged pores, and shallow post-acne scarring.
            \item Rough, uneven skin texture (hyperkeratosis).
        \end{itemize}
    \item \textbf{Expected post-procedure skin features:}
        \begin{itemize}
            \item Mild, transient erythema (redness) and a sensation of warmth resembling a mild sunburn, which typically resolves within 2 to 24 hours.
            \item Minimal localized oedema (swelling), giving the skin a temporarily plumped appearance.
            \item \textbf{Mild desquamation:} Fine flaking or peeling of the skin may occur 2 to 3 days post-treatment as the exfoliated layers shed.
        \end{itemize}
\end{itemize}

\section*{Differential Diagnosis}
Microdermabrasion must be distinguished from other skin resurfacing and rejuvenation techniques that penetrate different layers of the skin:
\begin{itemize}
    \item \textbf{Dermabrasion:} A highly invasive surgical procedure performed under local or general anaesthesia. Dermabrasion uses a rapidly rotating wire brush or diamond burr to remove the entire epidermis and penetrate the reticular dermis, carrying a significant risk of scarring, permanent pigment changes, and prolonged recovery.
    \item \textbf{Chemical peels:} Resurfacing achieved via the application of chemical acids (e.g., salicylic, glycolic, or trichloroacetic acid) that cause controlled chemical necrosis and shedding of the skin, rather than mechanical exfoliation.
    \item \textbf{Fractional laser resurfacing:} The use of laser energy to create microscopic zones of thermal damage in the skin to stimulate deep collagen remodeling, suitable for deep wrinkles but involving greater downtime and risk.
    \item \textbf{Microneedling (Collagen Induction Therapy):} The mechanical creation of thousands of microscopic puncture wounds in the dermis using sterile needles, rather than superficial surface exfoliation.
\end{itemize}

\section*{When to Seek Medical Care}
Patients should seek medical advice from their treating clinician or a general practitioner if post-procedure skin changes do not resolve as expected, or if symptoms suggestive of an adverse reaction develop.

\textbf{Contact your local emergency services for:}
\begin{itemize}
    \item Severe, immediate systemic allergic reactions (anaphylaxis) to topical preps, cleansers, or numbing creams used during the session, presenting with swelling of the face, lips, or tongue, and difficulty breathing.
    \item High fever, severe spreading pain, and rapid swelling of the treated skin (suggesting spreading bacterial cellulitis).
    \item Deep skin ulceration or necrosis, which is extremely rare and only occurs with inappropriate use of non-medical devices.
\end{itemize}

\section*{Diagnosis and Investigations}
No laboratory investigations are required. Assessment involves a clinical evaluation by a qualified practitioner prior to the procedure.
\begin{itemize}
    \item \textbf{Skin assessment and Fitzpatrick typing:} Determining the patient's skin type and documenting baseline photographs to monitor treatment progress.
    \item \textbf{Screening for contraindications:} Identifying absolute or relative contraindications, which include:
        \begin{itemize}
            \item Active inflammatory acne (pustules or cysts) that could be ruptured and spread by mechanical friction.
            \item \textbf{Active skin infections:} Such as herpes simplex virus (cold sores), impetigo, or molluscum contagiosum in the treatment area.
            \item Vasculopathic or inflammatory conditions: Severe rosacea, active eczema, or psoriasis, which can be flared by mechanical stimulation.
            \item Recent medication use: The use of oral isotretinoin (Roaccutane) within the preceding 6 to 12 months, which impairs skin healing and increases the risk of scarring.
        \end{itemize}
\end{itemize}

\section*{Management}
The procedure and post-treatment phase must be carefully managed to ensure safety and optimal outcomes.
\begin{itemize}
    \item \textbf{Procedure protocol:}
        \begin{itemize}
            \item Cleansing the skin thoroughly to remove lipids and cosmetic residue.
            \item \textbf{Eye protection:} Placing protective goggles or pads over the patient's eyes to prevent corneal injury from stray crystals or dust.
            \item Exfoliation: Systematically moving the handpiece across the treatment area using gentle, sweeping strokes, adjusting the vacuum pressure according to skin thickness (reducing pressure over the delicate periorbital area).
        \end{itemize}
    \item \textbf{Post-procedure skin care:}
        \begin{itemize}
            \item \textbf{Soothing topicals:} Application of a gentle, fragrance-free moisturiser containing ceramides or hyaluronic acid to support epidermal barrier recovery.
            \item \textbf{Strict photoprotection:} Mandatory daily application of a broad-spectrum SPF 50+ sunscreen, as the exfoliated epidermis is highly susceptible to ultraviolet (UV) radiation damage and hyperpigmentation.
            \item \textbf{Skincare modification:} Avoidance of active topical ingredients (such as retinoids, salicylic acid, glycolic acid, or benzoyl peroxide) and physical scrubs for at least 3 to 5 days post-procedure to prevent severe contact irritation.
        \end{itemize}
\end{itemize}

\section*{Complications}
While microdermabrasion is low-risk, adverse events can occur, particularly if performed by untrained operators.
\begin{itemize}
    \item Prolonged erythema or mild localized oedema lasting more than 48 hours.
    \item Petechiae (micro-bruising) or purpura, which can occur if the vacuum pressure is set too high or if the patient has fragile capillaries.
    \item Reactivation of herpes simplex virus (cold sores) triggered by mechanical trauma to the perioral area.
    \item Post-inflammatory hyperpigmentation, rare but possible if the patient is exposed to direct sunlight without sun protection immediately following the session.
    \item Superficial bacterial or fungal infections resulting from poorly sanitised device tips or handpieces.
    \item Ocular irritation or corneal abrasion if crystals enter the eyes due to inadequate protective eyewear.
\end{itemize}

\section*{Prognosis and Outcomes}
The prognosis for skin rejuvenation following microdermabrasion is excellent. Recovery is immediate, and patients can resume normal daily activities without downtime. While a single treatment provides temporary skin smoothing and brightness, a clinical series of 4 to 6 sessions, spaced 2 to 4 weeks apart, is required to achieve noticeable improvements in fine lines, hyperpigmentation, and acne scars. Maintenance treatments are recommended every 1 to 2 months thereafter. The procedure does not correct deep wrinkles, structural sagging, or deep boxcar/icepick acne scars, which require more invasive dermatological therapies.

\section*{Prevention and Self-Care}
To ensure safety and maximize the benefits of microdermabrasion, patients should adhere to the following self-care guidelines:
\begin{itemize}
    \item Ensure the procedure is performed by a licensed, reputable dermal clinician, aesthetic nurse, or dermatologist who maintains strict hygiene standards.
    \item Avoid direct sun exposure and tanning beds for at least one week prior to and two weeks after each treatment.
    \item Keep the skin well-hydrated post-procedure using gentle, bland moisturisers.
    \item Do not pick, scratch, or peel any flaking skin post-treatment to avoid scarring and secondary infection.
    \item Discontinue any potentially irritating skincare products several days before the scheduled session.
\end{itemize}

\section*{Sources and Further Reading}
The following reputable organisations provide further information and safety advice on microdermabrasion and cosmetic procedures:
\begin{itemize}
    \item Australasian College of Dermatologists. \textit{Patient information on skin rejuvenation and cosmetic procedures.} https://www.dermcoll.edu.au/
    \item Healthdirect Australia. \textit{Trusted advice on non-surgical cosmetic procedures and choosing a practitioner.} https://www.healthdirect.gov.au/
    \item DermNet. \textit{Clinical articles detailing microdermabrasion indications, contraindications, and complications.} https://dermnetnz.org/
    \item American Society of Plastic Surgeons (ASPS). \textit{Aesthetic guide to microdermabrasion procedure and recovery.} https://www.plasticsurgery.org/
    \item Mayo Clinic. \textit{Overview of cosmetic skin resurfacing techniques.} https://www.mayoclinic.org/
\end{itemize}
